\documentclass[12pt]{article}
\usepackage[T1]{fontenc}
\usepackage[utf8]{inputenc}
\usepackage{lmodern}
\usepackage{textcomp}
\usepackage{enumitem}
\usepackage{geometry}
\geometry{margin=1in}
\begin{document}
\begin{center}
Answer Key - Sociology - Graduate Level Assessment 2 \
\small (Answers are ordered exactly as in the question paper.)
\end{center}

\section*{Multiple Choice}
\begin{enumerate}[label=\arabic*.]
\item \textbf{Answer:} C
\\ \textit{Options:} A) Ethnocentrism; B) Stereotyping; C) Scapegoating; D) Institutionalization
\\ \textit{Note:} Scapegoating refers to the social phenomenon where a specific group is unjustly blamed for the wider societal issues, often as a means for others to deflect responsibility or avoid confronting the complex causes of those problems.
\item \textbf{Answer:} C
\\ \textit{Options:} A) the formation of an attachment bond between an infant and its carer; B) a tendency of social theorists to explain everything in terms of social causes; C) the process of becoming part of a society by learning its norms and values; D) the historical process by which societies change from traditional to modern
\\ \textit{Note:} The correct answer defines socialization as the essential process through which individuals internalize the norms and values of their society, facilitating their integration and function within that social framework.
\item \textbf{Answer:} A
\\ \textit{Options:} A) a subtler form of prejudice, masked by nationalist pride; B) a post-modern deconstruction of racist ideas to reveal their lack of depth; C) racist practices found in newly emerging areas of social life, such as cyberspace; D) an anti-fascist movement which challenges nationalist politics
\\ \textit{Note:} The 'new racism' is characterized by its subtlety and the way it is often expressed through nationalist pride, rather than overtly discriminatory language or actions, reflecting a shift in how prejudice manifests in contemporary society.
\item \textbf{Answer:} A
\\ \textit{Options:} A) transnational corporations; B) multi-media empires; C) ownership concentrated within one medium; D) government-owned companies
\\ \textit{Note:} Walt Disney, Sony, and Time Warner are considered transnational corporations because they operate and have significant business activities in multiple countries, enabling them to influence global markets and cultures.
\item \textbf{Answer:} B
\\ \textit{Options:} A) impersonal rules; B) extensive paperwork; C) hierarchy of officials; D) Salaries
\\ \textit{Note:} Extensive paperwork is not part of Weber's ideal type of bureaucracy because he emphasized characteristics such as hierarchical organization, division of labor, and formal rules, rather than the volume of documentation involved.
\end{enumerate}

\section*{True / False}
\begin{enumerate}[label=\arabic*.]
\setcounter{enumi}{5}
\item \textbf{Answer:} True
\\ \textit{Options:} A) True; B) False
\\ \textit{Note:} A household account book is considered a personal document with closed access because it contains sensitive financial information that reflects the private economic activities and decisions of the household, thereby necessitating confidentiality to protect individual privacy and autonomy.
\item \textbf{Answer:} False
\\ \textit{Options:} A) True; B) False
\\ \textit{Note:} Weber argued that the 'spirit of capitalism' was influenced by Protestant ethics, particularly Calvinism, which promoted a work ethic and rationalization of labor rather than a rejection of traditional labor practices.
\item \textbf{Answer:} True
\\ \textit{Options:} A) True; B) False
\\ \textit{Note:} The New Right primarily emphasized traditional family structures and moral conservatism, often overlooking cohabitation as a relevant factor in their critiques of family values.
\item \textbf{Answer:} True
\\ \textit{Options:} A) True; B) False
\\ \textit{Note:} Secondary deviation refers to the process by which an individual's deviant behavior is amplified through societal reactions, including punishment and labeling, leading them to adopt a deviant identity and engage in further deviance as a response to societal stigma.
\item \textbf{Answer:} True
\\ \textit{Options:} A) True; B) False
\\ \textit{Note:} Berger (1963) employs the puppet theatre metaphor to demonstrate how individuals are socialized into roles and scripts that shape their understanding of reality, highlighting the constructed nature of social interactions and institutions.
\end{enumerate}

\section*{Long Form Response}
\begin{enumerate}[label=\arabic*.]
\setcounter{enumi}{10}
\item \textbf{Answer:} Current research findings indicate that most men married to working women do not share housework and childcare obligations equally with their spouses, highlighting a persistent imbalance in domestic responsibilities. This division of labor reinforces traditional gender roles, where women remain primary caregivers and homemakers despite their professional commitments. Such dynamics not only perpetuate gender inequality but also affect women's career advancement and overall well-being. As both partners juggle work responsibilities, the disproportionate allocation of domestic tasks can lead to increased stress and conflict within marriages, ultimately impacting family dynamics and societal perceptions of gender roles. Addressing these disparities is crucial for fostering equality in both the workplace and home, thereby supporting a more equitable distribution of responsibilities among couples.
\\ \textit{Note:} The answer correctly identifies the unequal distribution of housework and childcare responsibilities among married couples where both partners work, illustrating how this imbalance perpetuates traditional gender roles and negatively impacts women's career progression and well-being.
\item \textbf{Answer:} In sociology, an aggregate refers to a collection of individuals who come together in a specific location but do not share a common identity or interact meaningfully. The scenario of seven people waiting silently at a bus stop exemplifies this concept, as these individuals are physically present in the same space but lack any significant social connection or purpose beyond awaiting transportation. This situation highlights the characteristics of an aggregate, such as the absence of structured relationships or shared goals, as the individuals may not engage in conversation or form a cohesive group. Instead, their proximity is purely coincidental, underscoring the distinction between aggregates and more integrated social entities like groups or communities. Thus, the bus stop scenario effectively illustrates how aggregates are defined by their lack of interaction and shared identity, despite their simultaneous presence.
\\ \textit{Note:} correct because it clearly defines an aggregate as a collection of individuals who share a physical space without meaningful interaction or shared identity, effectively illustrated by the example of seven people silently waiting at a bus stop.
\end{enumerate}

\vfill
\noindent\textit{End of Answer Key}
\end{document}